\section{A unitary with a large fixed space}\label{sec:degree}

The next lemma supplies the matrix that will mix copies of a word disk.
Its conclusion holds for an arbitrary fixed unitary matrix; no continuous
choice of the mixing matrix is needed.
The compact-unitary degree strategy goes back to
Gerstenhaber--Rothaus \cite[Theorem 1]{GerstenhaberRothaus1962};
see also \cite[Section 1.1]{KlyachkoThom2017} for removing coefficients
by homotopy. Here the degree domain records a fixed plane as well as
a unitary, so that it has the dimension required for the intersection.

\begin{lemma}\label{lem:fixed}
Let $s_1,\ldots,s_n\in\{1,-1\}$ satisfy
$\varepsilon:=\sum_i s_i\in\{1,-1\}$, let $k\geq1$ be an integer,
and put $m=nk$.
For every $P\in\U(m)$ there is an $X\in\U(k)$ such that
\[
 F(X)=\operatorname{diag}(X^{s_1},\ldots,X^{s_n})P
\]
fixes a complex subspace of dimension at least $k$.
\end{lemma}

\begin{proof}
We use degree modulo two to force an intersection with the matrices
that fix a $k$-plane.  Keeping the plane as part of the data makes that
incidence space a smooth manifold.

\medskip\noindent
\emph{A convenient homotopy.}
First remove $P$ along a smooth path from $P$ to $I_m$ in $\U(m)$,
obtained by diagonalizing this fixed matrix and interpolating its angles.
All homotopies below can be reparametrized to be constant near their
endpoints, so their concatenation is smooth.
Next pair positive and negative blocks, leaving one block of sign
$\varepsilon$.  Each pair can be cancelled by the following explicit
homotopy.  On its two coordinate $k$-planes, set
\[
 R_u=\begin{pmatrix}
 \cos(\pi u/2)I_k&-\sin(\pi u/2)I_k\\
 \sin(\pi u/2)I_k&\cos(\pi u/2)I_k
 \end{pmatrix},\qquad 0\leq u\leq1.
\]
Then
\[
 \operatorname{diag}(X,I_k)R_u
 \operatorname{diag}(I_k,X^{-1})R_u^{-1}
\]
runs from $\operatorname{diag}(X,X^{-1})$ to $I_{2k}$ and depends
smoothly on $X$.  Apply this homotopy on the disjoint paired block
sums, and then move the resulting constant identity blocks to $-I$.
Thus, in a fixed orthogonal splitting
$\C^m=E_0\oplus E_0^\perp$ with $\dim_{\C}E_0=k$, the map $F$ is
homotopic to
\[
 F_0(X)=X^\varepsilon\oplus(-I_r),\qquad r=m-k.
\]

\medskip\noindent
\emph{The incidence manifold.}
Write $\Gr_k(\C^m)$ for the Grassmannian of complex $k$-planes and set
\[
 \mathcal R=\{(W,E)\in\U(m)\times\Gr_k(\C^m):
                 W|_E=I_E\}.
\]
A unitary fixing $E$ preserves $E^\perp$, where it is arbitrary.
Consequently $\mathcal R$ is a smooth bundle over $\Gr_k(\C^m)$
with fiber $\U(r)$.  Local orthonormal frames of $E^\perp$ give its
local product charts.  The condition $(W-I_m)p_E=0$, where $p_E$ is
the orthogonal projection onto $E$, also shows that $\mathcal R$ is
closed in a compact product.  Its base and fiber are connected.
Thus it is a compact connected manifold without boundary, of real dimension
\[
 2kr+r^2=m^2-k^2.
\]
The product $M=\U(k)\times\mathcal R$ and the target $\U(m)$ are
thus connected closed manifolds of the same dimension $m^2$.

Consider the smooth map
\[
 G_0:M\longrightarrow\U(m),\qquad
 G_0(X,W,E)=F_0(X)W^{-1}.
\]
Its only preimage of $I_m$ is
\[
 (I_k,W_0,E_0),\qquad W_0=I_k\oplus(-I_r).
\]
Indeed, $G_0(X,W,E)=I_m$ implies $W=F_0(X)$.  Every fixed vector
of this matrix lies in $E_0$, so its fixed $k$-plane $E$ must equal
$E_0$.  It follows that $X^\varepsilon=I_k$, hence $X=I_k$.

\medskip\noindent
\emph{Regularity at the unique preimage.}
For $j\geq0$, let $\mathfrak u(j)$ denote the real vector space of
skew-Hermitian $j\times j$ matrices.  Tangent coordinates at this
preimage are
\[
 A\in\mathfrak u(k),\qquad B\in\mathfrak u(r),\qquad
 Z\in M_{r\times k}(\C).
\]
Put
\[
 K_Z=\begin{pmatrix}0&-Z^*\\Z&0\end{pmatrix}.
\]
Local coordinates are given by
\[
 X=e^A,\qquad E=e^{K_Z}E_0,\qquad
 W=e^{K_Z}\operatorname{diag}(I_k,-e^B)e^{-K_Z}.
\]
Here $Z$ gives the usual Grassmannian tangent directions
$\operatorname{Hom}_{\C}(E_0,E_0^\perp)$, and $B$ gives all fiber
directions.  At zero these formulas give
\[
 \dot F_0=\begin{pmatrix}\varepsilon A&0\\0&0\end{pmatrix},
 \qquad
 \dot W=[K_Z,W_0]+\begin{pmatrix}0&0\\0&-B\end{pmatrix}
       =\begin{pmatrix}0&2Z^*\\2Z&-B\end{pmatrix}.
\]
Since $F_0=W=W_0=W_0^{-1}$ at the preimage,
\begin{equation}\label{eq:degree-derivative}
 dG_0(A,B,Z)=\dot F_0W_0-\dot WW_0
   =\begin{pmatrix}\varepsilon A&2Z^*\\-2Z&-B\end{pmatrix}.
\end{equation}
Every skew-Hermitian $m\times m$ matrix has this form uniquely.
Thus the derivative is an isomorphism and $I_m$ is a regular value.
If $r=0$, the Grassmannian and $\mathcal R$ are points, and the
calculation reduces to $A\mapsto\varepsilon A$.  The same formulas
also apply when $k=1$.

\medskip\noindent
\emph{The degree conclusion.}
For smooth maps between connected closed manifolds of the same dimension,
the parity of a regular fiber is independent of the regular value and
invariant under smooth homotopy. This is the degree modulo two;
see \cite[Section 4, pp.\ 20--25]{Milnor1965}. No orientations are needed.
The unique regular preimage above therefore gives $\deg_2(G_0)=1$.

The homotopy from $F_0$ to $F$ induces a homotopy from $G_0$ to
\[
 G(X,W,E)=F(X)W^{-1}.
\]
Hence $G$ also has degree one.  It cannot miss $I_m$, since a map
with no preimages of that point has degree zero.  At a preimage
we have $F(X)=W$, and $W$ fixes $E$.  This proves the lemma.
\end{proof}
